\chapter{Faithful CTL}\label{ch:faithful}

This chapter turns a protocol into CTL formulas that say exactly what the protocol
allows: in each state the conversation can be in, each next state is possible and
nothing else is. It proves that a system satisfies the formulas if and only if it is
bisimilar to the protocol, which by Browne's result (Section~\ref{thm:bcg}) is the most any CTL can say.

\section{Reading CTL}

CTL talks about a conversation as a tree of possible futures. Standing at one moment of
the conversation, the next message could be one of several; after each of those, again
one of several; and so on. A CTL formula is a statement made at one moment about that
tree. This chapter needs only six pieces, and each has a plain reading.

\begin{center}
\begin{tabular}{@{}lll@{}}
\toprule
written & CTL & read it as \\
\midrule
$\on{GetImage}$ & & on delivery of GetImage \\
$\term$ & & the conversation may stop here \\
$\someNext\varphi$ & $\mathsf{EX}$ & \emph{some} possible next step leads to a moment where $\varphi$ \\
$\allNext\varphi$ & $\mathsf{AX}$ & \emph{every} possible next step leads to a moment where $\varphi$ \\
$\always\varphi$ & $\mathsf{AG}$ & $\varphi$ holds at every moment the conversation can ever reach \\
$\varphi\rightarrow\psi$ & & whenever $\varphi$, also $\psi$ \\
\bottomrule
\end{tabular}
\end{center}

$\wedge$ is ``and'', $\vee$ is ``or'', $\neg$ is ``not''. Chapter~\ref{ch:checker} gives
the same definitions formally; this table is all the next section uses.

\section{The idea, on one request}

Take Images (Chapter~\ref{ch:global}): a browser sends GetImage, and the server answers
with an Image or a NotFound, and that is the whole conversation. Stand at the moment
just after GetImage. Three things are true there, and each is one line of CTL.

\begin{center}
\begin{tabular}{@{}lp{.36\textwidth}@{}}
\toprule
CTL & in words \\
\midrule
$\someNext(\on{Image}\wedge\term)$ & an Image \emph{can} come next, and then the conversation may stop \\
$\someNext(\on{NotFound}\wedge\term)$ & a NotFound \emph{can} come next, and then it may stop \\
$\allNext\big((\on{Image}\wedge\term)\vee(\on{NotFound}\wedge\term)\big)$ & \emph{whatever} comes next is one of those two; nothing else can \\
\bottomrule
\end{tabular}
\end{center}

The first two lines rule out a server that never sends one of the answers (say, one
that cannot say 404). The third rules out a server that sends anything else (say, a
redirect, or nothing at all). Together they say: after the request, exactly these two
answers are possible.

That statement is about one moment. To make it a rule about the whole conversation,
say \emph{when} it applies: ``at every moment at which the last message was GetImage
and the conversation may not stop yet, those three lines hold.''
\begin{align*}
&\always\Big(\on{GetImage}\wedge\neg\term\ \rightarrow\\
&\qquad\someNext(\on{Image}\wedge\term)\wedge\someNext(\on{NotFound}\wedge\term)\\
&\qquad{}\wedge\allNext\big((\on{Image}\wedge\term)\vee(\on{NotFound}\wedge\term)\big)\Big)
\end{align*}
Write one such rule for every moment the protocol can be in, add one saying where it
starts, and the rules together say \emph{everything} the protocol says: what may happen
next at every point, and what may not. The rest of this chapter proves that, and does
the one part that is harder than it looks. Here the moment ``just after GetImage'' was
easy to name: the last message says it. In a larger protocol the same last message can
be heard at two different points (an Ack after the first request and an Ack after the
second), and the rule must tell them apart by what can still happen. That is what
typestates and their names are for.

\begin{figure}[ht]
\centering
\AccessibleFigure{Typestate graph of Images: start then GetImage then either Image or NotFound each with a double border because the session may stop there and a loop for the stutter.}{\Typestates[.62\textwidth]{ctl-images}}
\caption{The four moments of Images, as LeanFM computes them. A double border marks a moment where the conversation may stop; its loop says nothing more happens.}
\label{fig:images-ts}
\end{figure}

\section{Typestates}

A \emph{typestate} is a state of the protocol as far as anyone can tell: two states are the
same typestate when they have the same label (the last message, and whether the conversation
may stop) and the same typestates can come next. LeanFM finds them by splitting groups until
nothing splits any more:

\begin{definition}[Refinement rounds]\label{def:rounds}
Start with one group per label, and split until nothing splits:
\begin{lstlisting}[language={},basicstyle=\ttfamily\footnotesize]
group = {s: label(s) for s in States}                  # round 0: last message, may stop
while True:
    split = {s: (group[s], frozenset(group[t] for t in next(s))) for s in States}
    if len(set(split.values())) == len(set(group.values())):
        break                                          # nothing split: done
    group = split
\end{lstlisting}
The groups left at the end are the \emph{typestates}; the number of rounds that split
something is the \emph{lookahead}.
\end{definition}

\begin{lemma}[Typestates are the same-moves classes]\label{lem:rounds}
The typestates are exactly the classes of states with the same moves (bisimilarity,
Section~\ref{def:bisim}).
\end{lemma}

\begin{proof}
When the loop stops, states in one group have the same label and step into the same groups,
so each move of one is matched by a move of the other into the same group: same moves. And
the loop never separates two states with the same moves, since they have the same label and
step into the same groups in every round.
\end{proof}

The lookahead is how far ahead one must look to tell typestates apart. At $0$ the label
alone decides: Images is like that (Figure~\ref{fig:images-ts}). The Retry loop of
Chapter~\ref{ch:constructs} is not: a Busy after the first Ask and a Busy after the third
look the same, and only what may come next (another Ask, or GiveUp) tells them apart.

Another way to read a typestate: it is the message, plus the bits it would have to carry
for the message to say what may come next. A name that labels one typestate carries nothing
more; a name that labels $n$ of them needs $\lceil \log_2 n \rceil$ bits, in a field, by
cause, or remembered by every participant. For Retry, from the same count \texttt{session\_ctl}
answers as \texttt{bits}:
\begin{center}
\begin{tabular}{@{}lrr@{}}
\toprule
Message & Typestates & Bits \\
\midrule
\texttt{start} & 1 & 0 \\
\texttt{Ask} & 3 & 2 \\
\texttt{Ok} & 1 & 0 \\
\texttt{Busy} & 3 & 2 \\
\texttt{GiveUp} & 1 & 0 \\
\bottomrule
\end{tabular}

\end{center}
Ask and Busy each need two bits: the try counter. Lookahead $4$ is what it costs to leave
those bits off the wire and recover them from what follows. Images needs none: its names are
its states.

\section{Naming a typestate}

Each typestate gets a formula that is true exactly at its states, its \emph{name}, built up
round by round like the groups:

\begin{definition}[Names]\label{def:chi}
A typestate's name at each round:
\begin{lstlisting}[language={},basicstyle=\ttfamily\footnotesize]
name(C, 0)   = (state M)  or  (state M stop)  or  start       ; its label
name(C, k+1) = (and name(C, 0)
                    (someNext name(D1, k)) ... (someNext name(Dn, k))
                    (allNext (or name(D1, k) ... name(Dn, k))))
               ; D1 ... Dn: the groups of round k that C's states step into
\end{lstlisting}
A typestate's name is its name at the last round.
\end{definition}

At lookahead $0$ a name is just a state term, \texttt{(state GetImage)}. With lookahead, a
name also says, as many steps deep as needed, what may come next. Two facts make the names
work:

\begin{lemma}[Each state satisfies its name]\label{lem:own}
Every state of the protocol satisfies the name of its own typestate.
\end{lemma}

\begin{proof}
Round by round: a state has its own label, and at each round its next states lie in the
groups $D_i$ it steps into (the $\mathsf{allNext}$ part) and it reaches each of them (the
$\mathsf{someNext}$ parts), since every state of a group steps into the same groups.
\end{proof}

\begin{lemma}[Names exclude each other]\label{lem:exclusive}
In any system, no state satisfies the names of two different typestates.
\end{lemma}

\begin{proof}
Round by round: two groups split at round $0$ have different labels, which conflict. Two
groups split later step into different groups, so one has a $\mathsf{someNext}$ the other's
$\mathsf{allNext}$ rules out.
\end{proof}

\section{The rules}

The rules are the transition table: where the conversation starts, and for each typestate
which typestates may come next.

\begin{definition}[Faithful rules]\label{def:rules}
The start, and one row per typestate:
\begin{lstlisting}[language={},basicstyle=\ttfamily\footnotesize]
name(start typestate)                                  ; it begins there
(follows name(C) name(D1) ... name(Dn))                ; for each typestate C,
                                                       ; D1 ... Dn the ones C steps into
\end{lstlisting}
With lookahead $0$ every name is a state term, so each \texttt{follows} is a
\texttt{transition} row, or \texttt{final} where a typestate steps only into itself.
\end{definition}

\begin{theorem}[Faithfulness]\label{thm:faithful}
A system satisfies the faithful rules exactly when it has the protocol's moves at every
state, no move more and no move less (it is bisimilar to the protocol).
\end{theorem}

\begin{proof}
If it has the same moves: the protocol satisfies its own rules (each state satisfies its
name, and its next states are exactly the typestates its row lists), and a system with the
same moves satisfies the same formulas (Browne's result, Section~\ref{thm:bcg}).

If it satisfies the rules: pair each reachable state of the system with the protocol
typestate whose name it satisfies, at most one by the exclusion of names (Section~\ref{lem:exclusive}); the start rule
pairs the two starts. Paired states have the same label, which every name fixes. Each move
the system makes lands in a typestate its row lists (the $\mathsf{allNext}$ part), so the
protocol has that move; each move the protocol makes is one the system can make (the
$\mathsf{someNext}$ parts). So paired states have the same moves, step after step: the
system is bisimilar to the protocol.
\end{proof}

\begin{corollary}[Nothing in CTL does better]\label{cor:best}
No set of CTL formulas tells more systems apart from the protocol than its faithful rules,
and every CTL property of the protocol holds of every system that satisfies them.
\end{corollary}

\begin{proof}
A system that satisfies the rules has the protocol's moves (Section~\ref{thm:faithful}), so
it satisfies exactly the protocol's formulas (Browne's result, Section~\ref{thm:bcg}). A system that does not
satisfy the rules fails a formula the protocol satisfies.
\end{proof}

\section{The rules of Images}

LeanFM writes the rules with \texttt{session\_ctl} and \texttt{faithful: "rules"}. For
Images, with lookahead $0$, every name is a label and an ending:

\begin{lstlisting}[language={},basicstyle=\ttfamily\footnotesize,breaklines=true]
start
  start

typestate 0
  transition(start, GetImage)

typestate 1
  transition(GetImage, Image stop, NotFound stop)

typestate 2
  final(Image)

typestate 3
  final(NotFound)
\end{lstlisting}


The row \texttt{(transition GetImage (Image stop) (NotFound stop))} is the sentence this
chapter began with. The two \texttt{final} rows are the ended typestates: once an answer has
arrived, nothing else happens.

\section{A larger protocol}

PlaceOrder has 10 checker nodes, 10 typestates, lookahead 0\unskip. Its typestate graph is again
computed rather than drawn.

\begin{figure}[ht]
\centering
\AccessibleFigure{Typestate graph of PlaceOrder: start CreateOrder then Quote or SoldOut; after Quote Pay Charge then Charged and Receipt or Declined and PaymentDeclined. Receipt PaymentDeclined and SoldOut have double borders and stutter loops.}{\Typestates{ctl-placeorder}}
\caption{The ten typestates of PlaceOrder.}
\label{fig:placeorder-ts}
\end{figure}

\begin{lstlisting}[language={},basicstyle=\ttfamily\footnotesize,breaklines=true]
start
  start

typestate 0: (start)
  transition(start, CreateOrder)

typestate 1: CreateOrder
  transition(CreateOrder, Quote, SoldOut stop)

typestate 2: Quote
  transition(Quote, Pay)

typestate 3: SoldOut (may stop)
  final(SoldOut)

typestate 4: Pay
  transition(Pay, Charge)

typestate 5: Charge
  transition(Charge, Charged, Declined)

typestate 6: Charged
  transition(Charged, Receipt stop)

typestate 7: Declined
  transition(Declined, PaymentDeclined stop)

typestate 8: Receipt (may stop)
  final(Receipt)

typestate 9: PaymentDeclined (may stop)
  final(PaymentDeclined)
\end{lstlisting}


\section{What the rules buy}

Because nothing in CTL does better than the rules (Section~\ref{cor:best}), every property one could ask of the protocol is a
consequence of eleven formulas. For example, every request is answered:
\[
\always\big(\on{CreateOrder} \rightarrow \inevitably(\on{Receipt}\vee\on{PaymentDeclined}\vee\on{SoldOut})\big)
\]
holds of PlaceOrder, so it holds of every system bisimilar to it, and need not be
stated separately. Properties that are \emph{not} consequences of the protocol, such as
a rule a project adds (``no refund before a charge''), are stated in the same logic and
checked against the same model. Chapter~\ref{ch:mcp} shows both kinds of question asked
of LeanFM.
